\subsection{Formula}

Let $D$ be a simply connected planar district with area $A(D)$ and perimeter $P(D)$.
The Polsby-Popper score is:
\begin{equation}
  \text{PP}(D) = \frac{4\pi A(D)}{P(D)^2}.
  \label{eq:pp}
\end{equation}

\begin{theorem}[PP is maximised by a disk]
  For any simply connected planar region $D$, $\text{PP}(D) \leq 1$,
  with equality if and only if $D$ is a disk.
\end{theorem}

\begin{proof}
  The isoperimetric inequality states $4\pi A \leq P^2$ for any simple closed curve,
  with equality if and only if the curve is a circle~\citep{osserman1978isoperimetric}.
  Dividing both sides by $P^2$ gives $\text{PP}(D) = 4\pi A/P^2 \leq 1$,
  with equality iff $D$ is a disk. \qed
\end{proof}

\subsection{Canonical Values}

We record PP for common geometric shapes.

\begin{table}[ht]
\centering
\caption{Polsby-Popper score for canonical shapes.}
\label{tab:canonical}
\begin{tabular}{lcc}
\toprule
Shape & PP & Approximate value \\
\midrule
Disk (circle) & $1$ & $1.000$ \\
Regular hexagon & $\pi\sqrt{3}/6$ & $0.907$ \\
Square         & $\pi/4$        & $0.785$ \\
Equilateral triangle & $\pi/(3\sqrt{3})$ & $0.605$ \\
$1 \times 3$ rectangle & $3\pi/16$ & $0.589$ \\
$1 \times 10$ rectangle & $10\pi/121$ & $0.260$ \\
\bottomrule
\end{tabular}
\end{table}

\subsection{Properties}

\begin{proposition}
  PP is scale-invariant: $\text{PP}(\lambda D) = \text{PP}(D)$ for any $\lambda > 0$.
\end{proposition}

\begin{proof}
  Scaling by $\lambda$ gives $A \mapsto \lambda^2 A$ and $P \mapsto \lambda P$,
  so PP $= 4\pi(\lambda^2 A)/(\lambda P)^2 = 4\pi A/P^2$. \qed
\end{proof}

\begin{proposition}
  PP is rotation-invariant and translation-invariant.
\end{proposition}

These invariances confirm PP is a geometric property of the shape,
independent of its location or orientation on the map.

\subsection{Relation to Schwartzberg}

The Schwartzberg score $S = P/(2\sqrt{\pi A})$ satisfies $S = 1/\sqrt{\text{PP}}$
exactly (K.4, Theorem~1). This means PP and S carry identical geometric information
and can be converted losslessly: a district with PP $= 0.22$ has
$S = 1/\sqrt{0.22} \approx 2.13$.
